\documentclass[11pt]{report}

% -------------------------------------------------------
% Page layout
% -------------------------------------------------------
\usepackage[pdftex]{geometry}
\geometry{
  a4paper,
  left=2.5cm,
  right=2.5cm,
  top=1.5cm,
  bottom=2.0cm,
  twoside
}

% -------------------------------------------------------
% Encoding, fonts, and language support
% -------------------------------------------------------
\usepackage[utf8]{inputenc}
\usepackage[T1]{fontenc}
\usepackage{lmodern}
\usepackage[main=english,spanish,brazilian]{babel}
\usepackage{textcomp}

% -------------------------------------------------------
% Math and symbols
% -------------------------------------------------------
\usepackage{amsmath}
\usepackage{amssymb}
\usepackage{enumerate}
\usepackage{bm}
\usepackage{bbm}
\usepackage{siunitx}

% -------------------------------------------------------
% Figures and tables
% -------------------------------------------------------
\usepackage{graphicx}
\graphicspath{{./Comments/}}
\usepackage{subfigure}
\usepackage{float}

% -------------------------------------------------------
% Colors and boxes
% -------------------------------------------------------
\usepackage{xcolor}
\usepackage[many]{tcolorbox}
\tcbuselibrary{breakable}
% -------------------------------------------------------
% References and links
% -------------------------------------------------------
\usepackage{natbib}
\usepackage{hyperref}

% -------------------------------------------------------
% Unicode fallback declarations for pdfLaTeX
% -------------------------------------------------------
\DeclareUnicodeCharacter{00E1}{\'a} % �
\DeclareUnicodeCharacter{00E9}{\'e} % �
\DeclareUnicodeCharacter{00ED}{\'i} % �
\DeclareUnicodeCharacter{00F3}{\'o} % �
\DeclareUnicodeCharacter{00FA}{\'u} % �
\DeclareUnicodeCharacter{00C1}{\'A} % �
\DeclareUnicodeCharacter{00C9}{\'E} % �
\DeclareUnicodeCharacter{00CD}{\'I} % �
\DeclareUnicodeCharacter{00D3}{\'O} % �
\DeclareUnicodeCharacter{00DA}{\'U} % �
\DeclareUnicodeCharacter{00F1}{\~n} % �
\DeclareUnicodeCharacter{00D1}{\~N} % �
\DeclareUnicodeCharacter{00E3}{\~a} % �
\DeclareUnicodeCharacter{00C3}{\~A} % �
\DeclareUnicodeCharacter{00E2}{\^a} % �
\DeclareUnicodeCharacter{00EA}{\^e} % �
\DeclareUnicodeCharacter{00F4}{\^o} % �
\DeclareUnicodeCharacter{00E7}{\c{c}} % �
\DeclareUnicodeCharacter{00C7}{\c{C}} % �
\DeclareUnicodeCharacter{00A0}{~} % non-breaking space
\DeclareUnicodeCharacter{2013}{--} % en dash
\DeclareUnicodeCharacter{2014}{---} % em dash
\DeclareUnicodeCharacter{2018}{`} % left single quote
\DeclareUnicodeCharacter{2019}{'} % right single quote
\DeclareUnicodeCharacter{201C}{``} % left double quote
\DeclareUnicodeCharacter{201D}{''} % right double quote

% -------------------------------------------------------
% Box styles
% -------------------------------------------------------
\newtcolorbox{reviewbox}[1]{
  colback = white!5!white,
  colframe = purple!75!black,
  fonttitle=\bfseries,
  title=#1,
  breakable,
  enhanced,
  before skip=10pt,
  after skip=10pt
}

\newtcolorbox{responsebox}[1]{
  colback = white!5!white,
  colframe = white!75!black,
  fonttitle=\bfseries,
  title=#1,
  breakable,
  enhanced,
  before skip=10pt,
  after skip=10pt
}

\newcommand{\assigned}[1]{\textcolor{blue}{\textbf{#1}}}

% -------------------------------------------------------
% Document
% -------------------------------------------------------
\begin{document}

\begin{center}
\large{\textbf{Response Letter to Reviewers}}

\vglue 0.3cm

\huge{
Adaptive Model-Free Tsallis Entropy Estimation for Detecting Non-Fully-Developed Speckle\\
(TGRS-2025-10064.R1)
}
\end{center}

\begin{center}
\textbf{
Janeth Alpala, Alejandro C.\ Frery, Paolo Gamba, Abra\~ao D. C. Nascimento, Andrea A. Rey
}
\end{center}

\date{ }


\noindent Dear Editor,

\bigskip

\noindent We are grateful for the careful evaluation of our revised manuscript and for the valuable comments provided in this round of review. We have addressed the remaining comments raised by Reviewer~2 and revised the manuscript accordingly. The corresponding changes have been highlighted in the revised version. 


\medskip

%%%%%%%%%%%%%%%%%%%%%%%%%%%%%%%%%%%%%%%%%%%%%%%%%%%%%%%%%%%%
\section*{Reviewer 1}

\begin{reviewbox}{Comment}

Accept.

\end{reviewbox}

\begin{responsebox}{Response}

We sincerely thank Reviewer~1 for the positive evaluation and acceptance of the revised manuscript.

\end{responsebox}

%%%%%%%%%%%%%%%%%%%%%%%%%%%%%%%%%%%%%%%%%%%%%%%%%%%%%%%%%%%%
\section*{Reviewer 2}

\begin{reviewbox}{Comment 1}

\textbf{Calibration and interpretation of the adaptive p-values}

Equation~(15) defines the adaptive z-score using $\sigma_{\mathrm{obs}}$, described as ``the empirical standard deviation of the test statistic computed over all the pixels under $H_0$ using adaptive windows.'' The following sentence then says that this scale is ``better matched to the variability of the observed SAR data.'' Since $\sigma_{\mathrm{obs}}$ is obtained from an $H_0$ Monte Carlo simulation rather than from each observed scene, I suggest clarifying this wording and the inferential role of the $H_0$-derived pooled scale used in Eq.~(15).
\vspace{1em}

It would also be helpful to state whether the adaptive calibration in Table~II and Eq.~(15) was recomputed for each $(L,\eta)$ pair used in the experiments. Section~III-G reports the adaptive null diagnostic for $\Gamma_{\mathrm{SAR}}(10,16)$ with $\eta = 3$, whereas the real-scene experiments use $L \in \{4,12,16,24\}$ and $\eta \in \{1,2,3,4\}$. Because the adaptive window selection depends on both $L$ and $\eta$, the null distribution of the adaptive statistic, and hence the calibration of the reported p-values, may differ across these settings if a single calibration is reused.

\vspace{1em}
A compact diagnostic of the final adaptive p-values under $H_0$, such as empirical rejection rates at 1\%, 5\%, and 10\%, or a check against a Uniform$(0,1)$ distribution, would strengthen the calibration argument. Such a diagnostic would be particularly important if a single calibration was reused across $L$ or $\eta$; if per-setting calibration was performed, a concise description of the calibration protocol and settings would likely be sufficient. Relatedly, the statement that moderate deviations in $L$ did not significantly affect the test outcome could be tied to the supporting sensitivity checks if those were performed.

\end{reviewbox}

\begin{responsebox}{Response}

We thank Reviewer 2 for this observation. We agree that the previous wording could lead to ambiguity regarding the origin of the scale used in Eq.~(15).

In the revised manuscript, we clarified the role of the location and scale terms in the adaptive $p$-value computation.
The simulation under $\mathcal H_0$ is used to estimate the window-dependent location, represented by the mean $\mu_{\mathcal H_0}(W)$, for each admissible window size $W\in\{5,7,9,11\}$.

We also clarified why the window-dependent scale, represented by $\sigma_{\mathcal H_0}(W)$, was not used directly.
Although this quantity was estimated in Table~II and evaluated as an alternative, using the scale obtained from the simulated image under $\mathcal H_0$ to standardize real SAR scenes led to excessive rejections in the $p$-value maps.
This indicated that such a scale was too small for the variability observed in real data.
Therefore, we used $\sigma_{\mathrm{obs}}$, the empirical standard deviation computed from the map of observed test-statistic values for the image under analysis.
This produced more stable $p$-value maps and better agreement with the ROI-based evaluation.

Following the reviewer's suggestion, we added a compact diagnostic of the adaptive $p$-values under $\mathcal H_0$.
Specifically, we added Fig.~7, which shows the histogram of the resulting $p$-values for the simulated homogeneous image.
The histogram is approximately flat and is consistent with the Uniform$(0,1)$ reference.
We also added Table~III, which reports the empirical rejection rates at the nominal levels $1\%$, $5\%$, and $10\%$.
These results support the behavior of the proposed adaptive $p$-value computation.

Regarding the dependence on $L$ and $\eta$, we clarified that the quantities estimated under $\mathcal H_0$ must correspond to the same $(L,\eta)$ pair used to process the image.
The number of looks $L$ is determined by each SAR scene, whereas $\eta$ is the adaptive-window parameter evaluated in the analysis.
Therefore, the values reported in Table~II should not be interpreted as universal values for all scenes or all choices of $\eta$.

For the Shannon, R\'enyi, and Tsallis benchmark, we clarified that the comparison was performed only on the Dublin scene, with $L=16$ and $\eta=3$.
For each entropy-based test, the corresponding quantities under $\mathcal H_0$ and the observed scale $\sigma_{\mathrm{obs}}$ were computed from its own test-statistic values.
The Dublin scene provides a clear separation between homogeneous water areas and heterogeneous urban structures, allowing a controlled ROI-based comparison among the three entropy-based tests.

These changes were incorporated in Section~III-G and Section~IV-E of the revised manuscript.

\end{responsebox}






\vspace{1em}

\begin{reviewbox}{Comment 2}

\textbf{Algebraic presentation of Eq. (8)}

The decomposition in Eq.~(8), $T_\lambda(G^0_I) = T_\lambda(\Gamma_{\mathrm{SAR}}) + \Delta_\alpha$, does not appear to be algebraically consistent with Eq.~(7) and the Appendix as displayed. If $Q = \lambda(L-1)+1$ and $R$ denotes the $\alpha$-dependent exponent, subtracting Eq.~(7) from the direct Appendix form gives an excess term proportional to $\{\lambda^{-Q} - \exp(R)\}$, not $\{1 - \exp(R)\}$. Equivalently, if the $-Q \ln \lambda$ term is absorbed into the exponential prefactor, the $\alpha$-dependent bracket changes as well.

The Appendix derivation and the R implementation appear to use the direct form $T = (1-I)/(\lambda-1)$. This appears to be purely a presentation issue rather than a problem with the experiments. I suggest revising Eq.~(8), or presenting $T_\lambda(G^0_I)$ directly, so that the displayed formula is consistent with Eq.~(7) and the Appendix.

\end{reviewbox}


\begin{responsebox}{Response}
We appreciate this detailed comment.  We agree with the reviewer that the displayed decomposition was not fully consistent with Eq.~(7) and with the direct expression derived in the Appendix.

After revising the formula, we identified that the previous version of Eq.~(8) inadvertently omitted one term inside the bracket. In the previous version, the bracket was written in the form
\begin{equation*}
\{1-\exp[\cdots]\}.
\end{equation*}
However, the correct bracket is
\begin{equation*}
\left\{
\exp\bigl[-\bigl(\lambda(L-1)+1\bigr)\ln\lambda\bigr]
-
\exp[\cdots]
\right\}.
\end{equation*}
The added term
\begin{equation*}
\exp\bigl[-\bigl(\lambda(L-1)+1\bigr)\ln\lambda\bigr]
\end{equation*}
is equivalent to
\begin{equation*}
\lambda^{-[\lambda(L-1)+1]},
\end{equation*}
which corresponds to the factor $\lambda^{-Q}$ mentioned by the reviewer, with
\begin{equation*}
Q=\lambda(L-1)+1.
\end{equation*}

Accordingly, we corrected Eq.~(8) in the revised manuscript. The corrected expression is
\begin{multline*}
T_\lambda\bigl(\mathcal{G}^0_I(\alpha, \mu, L)\bigr)
=
T_\lambda\bigl(\Gamma_{\mathrm{SAR}}(\mu, L)\bigr)
+
\frac{1}{\lambda-1}\,
\exp\Bigl[
(1-\lambda)\ln\mu
+(\lambda-1)\ln L
+\ln\Gamma\bigl(\lambda(L-1)+1\bigr)\\
-\lambda\ln\Gamma(L)
\Bigr]
\Bigl\{
\textcolor{blue}{\exp\bigl[-\bigl(\lambda(L-1)+1\bigr)\ln\lambda\bigr]}
-
\exp\bigl[
(1-\lambda)\ln(-\alpha-1)
+\lambda\ln\Gamma(L-\alpha)
-\lambda\ln\Gamma(-\alpha)\\
+\ln\Gamma\bigl(\lambda(1-\alpha)-1\bigr)
-\ln\Gamma\bigl(\lambda(L-\alpha)\bigr)
\bigr]
\Bigr\}.
\end{multline*}
The term added in blue is the missing factor from the $\Gamma_{\mathrm{SAR}}$ entropy expression in Eq.~(7). 
Its inclusion makes the decomposition in Eq.~(8) consistent with the direct expression derived in the Appendix and with the definition of the Tsallis entropy.

We also changed the notation from $\Delta_\alpha$ to $\Delta_\lambda(\alpha,\mu,L)$, because the additional term depends not only on the texture parameter $\alpha$, but also on the mean intensity $\mu$, the number of looks $L$, and the entropy order $\lambda$.

We emphasize that this correction affects only the algebraic presentation of the decomposition. 
The computational implementation and the numerical experiments are unchanged, because the R code uses the direct expression obtained from
\begin{equation*}
T_\lambda=\frac{1-I}{\lambda-1},
\end{equation*}
which is now explicitly displayed in the Appendix.

We used the exponential--logarithmic notation because it is consistent with the R implementation based on \texttt{log()} and \texttt{lgamma()}, and it improves numerical stability by avoiding the direct evaluation of large Gamma-function products.

The decomposition of $T_\lambda\bigl(\mathcal{G}^0_I(\alpha,\mu,L)\bigr)$ remains important for the theoretical interpretation of the proposed approach. It shows that the homogeneous $\Gamma_{\mathrm{SAR}}$ model is recovered as the limiting case of the $\mathcal{G}^0_I$ model when $\alpha\to-\infty$, while the additional term $\Delta_\lambda(\alpha,\mu,L)$ represents the texture-induced entropy contribution associated with heterogeneity.

\end{responsebox}




\vspace{1em}

\begin{reviewbox}{Comment 3}

\textbf{Scope of the revised claims}

The Section~IV-E benchmark is carried out on the Dublin scene only, so claims such as ``strongest overall performance'' should be scoped accordingly. Likewise, ``does not depend on sensors and resolutions'' overstates the evidence; ``behaved consistently across the tested configurations'' would be more precise.

\end{reviewbox}

\begin{responsebox}{Response}
We revised the manuscript to better delimit the scope of the claims and to align the wording with the evidence provided by the experiments. 
The comparative statements were reformulated more precisely, and the revised version now avoids overgeneralized expressions. 
The corresponding changes have been incorporated into the corrected manuscript and highlighted in the new version.

\end{responsebox}

\vspace{0.5cm}


\end{document}

